\chapter[How Early Empires Managed Difference]{How Early Empires Managed Differences}
\section{A Waterlogged Wooden Slip}
Pick something up: not gold or a monument, but wood, over twenty centimetres long, one or two wide, thinner than a ruler, bearing small brush-written lines. Over two millennia ago it was an official document in a northwestern Hunan county. Over two millennia later it emerged from an ancient well.

In 2002, excavators at Liye in Longshan County found over thirty thousand Qin wooden and bamboo documents, later called the Liye Qin slips. Not imperial edicts or histories, but Qianling County's daily files: household numbers, monthly grain receipts, an official's reprimand, document dispatch dates and deadlines. Repeated references to Qianling County, Dongting Commandery even corrected transmitted histories, which never recorded a Qin Dongting Commandery. Water-soaked wood revealed it.

Hold the slip and ask: western Hunan's mountain-enclosed Qianling lay on Qin's southwestern frontier, over a thousand kilometres of mountains and rivers from Xianyang near today's Xi'an. How did the emperor make this document appear here? Armies could not stand daily outside every county office. Voluntary initiative seems still less likely. How could a giant reach so far and so finely, into a minor clerk's account entry?

Deeper: that clerk was probably conquered Chu rather than Qin-born. Every character and rice measure followed conquerors' standards. Why no resistance? Was resistance possible? Or did he not think he served conquerors?

This mystery gives our subject: early empires managing difference.

\section{A Morning in Qianling County}
Enter a morning shortly after Qin unification, in Qianling beside the Youshui River.

Wuling mountain mist lingers after dawn; river damp clings to rammed-earth walls. Clerks already work in the county courtyard, rooms smelling of ink and damp wood. A middle-aged official kneels at a low desk before prepared slips, brush, and an official measuring vessel.

The first case concerns two farmers' disputed field boundary, complicated less by land than calculation. The elder uses inherited Chu acreage and harvest measures; the younger adopts official Qin units. Their figures cannot match. At the office door they argue over whose numbers count on this land.

The clerk knows the answer: imperial command makes standards one. Old Chu calculation is no longer an alternative, but wrong. He must translate dispute into new measures, write, file, report, using Qin-prescribed characters. His childhood Chu script becomes obsolete knowledge nobody teaches.

At noon, grain must move to the commandery. He carefully records quantity, date, handlers in rigid format, subject to checking and reprimand. Surviving slips genuinely preserve such penalties two thousand years later.

Evening brings familiar Chu bargaining and northern garrison accents outside. His pocket holds no Chu coin; vegetables require Qin banliang cash. He probably does not call his situation empire. The world, script, ruler changed; brush and ruler still earn his living.

Remember this morning. Its clerk's day contains every question that follows.

\section[Empires Face Differences, Not Maps]{Empires Face Differences, Not Merely Maps}
Set out the problem before answering.

One-colour territories around capitals make empires look like oversized states. Misunderstanding begins there. Early empires contain dozens or hundreds of peoples, languages, deities, land measures, silver scales, funerary rules. Conquest encloses them but cannot administer everyday life. Swords break walls, not routines.

Every empire therefore must answer:

\begin{itemize}
\item Money: what tribute and taxes, how much, and how transported to the centre?
\item People: who governs far-flung places? Insufficient insiders or potentially untrustworthy locals?
\item Routes: how do orders, armies, documents, grain cross journeys of months?
\item Rules: should measures, scripts, laws, money be unified, and how far?
\end{itemize}
Answers differ. Sargon's Akkad, founded around 2300 BCE in Mesopotamia and often called the first empire, conquered Sumerian cities, installed Akkadian governors, and replaced Sumerian with Akkadian for official documents. Egypt's New Kingdom, roughly sixteenth to eleventh centuries BCE, directly governed Nubia but kept Syrian and Palestinian princes as tributary clients sending sons as hostages. Assyria forcibly transplanted populations across its territory. Persia retained local gods, laws, rulers in exchange for tax and troops. Qin chose thorough standardisation: script, axle gauges, measures, dozens of commanderies and over a thousand counties, officials centrally appointed.

Circle this: \textbf{empire is not one fixed institution, but different solutions to shared problems}. The word sounds like a machine; better imagine an examination with radically different answers.

Can those answers be judged better or worse? Do common measures, script, money benefit ordinary people or erase origins? Does Xianyang's answer match Qianling's?

\section{Voices of Centre and Frontier}
Stand in different places and hear different imperial accounts.

\textbf{First: Xianyang's reasoning.}

Give the opposing case fully. Readers often reduce Qin standardisers to bullying, unfairly. Strengthen their arguments:

First, standards mean fairness. Warring States measures varied at borders, exposing merchants to extraction and taxpayers to adjustable rulers. A common measure blocks manipulation, hurting loophole exploiters. Second, standards lower everyone's costs: common axle gauges let carts use one road system; script lets a Qi scholar read documents as a Chu official; banliang money gives interregional trade a common reckoning. Third, deepest: centuries of destroyed towns and slaughter arose from a world divided into seven. Unified rules aim to prevent division's return. Ending war is their strongest card.

This case is substantial; modern versions still guide action worldwide.

\textbf{Second: the inscribed emperor.}

Persia's centre speaks from Behistun, western Iran, a vast inscription Darius I ordered around twenty-five centuries ago on a cliff a hundred metres high. Old Persian, Elamite, and Babylonian cuneiform proclaim the divinely protected king of kings repeatedly suppressing revolts, then warn posterity not to dismiss his true achievements as exaggeration.

Its stance celebrates conquest and obedience, the world under one king as achievement itself. Central chronicles count victories and treat differences as things to subdue and inventory. Yet three scripts acknowledge a multilingual realm even in royal proclamation.

\textbf{Third: weeping beside rivers.}

Change to conquered people. Assyria repeatedly deported conquered populations and replaced them with others. Late eighth-century BCE conquest of northern Israel brought removals and foreign resettlement in Samaria, as 2 Kings records. Imperial population allocation meant human uprooting.

Over a century later, successor hegemon Babylon captured Jerusalem and deported many Jews. Psalm 137 preserves their voice:

\begin{quote}
By the rivers of Babylon we sat down and wept when we remembered Zion.
\end{quote}
Zion means Jerusalem. No explicit accusation, only sitting and weeping, yet imperial accounts' relocated households acquire human shape. Central inscriptions count achievements; riverside poetry counts loss. One empire, two true voices.

\textbf{An easily misread counterexample.}

You may infer that empires always erase differences and homogenise everyone. Persia challenges that. After conquering Babylon, Cyrus permitted displaced peoples to return and restored captured gods to their temples. His cylinder portrays restoration rather than destruction; Ezra records permission for Jews to return and rebuild Jerusalem's temple.

New kindness? More complex: a different governing calculation retains local temples, laws, rulers, making differences serve empire. Elites continue local administration provided taxes, troops, loyalty follow. Erasure through deportation or direct county administration costs enormously; retention costs less. The lesson is not good Persia versus bad Assyria, but \textbf{genuinely different choices shaped by costs, ambition, worldview}, not simple virtue and vice.

\section[Thinking Bridge: An Imperial Check-Up]{A Thinking Bridge: Give an Empire a Check-Up}
How can one analyse a continental giant containing hundreds of peoples? Historians offer a portable scalpel: four diagnostic reports instead of memorised dates, useful later for states, companies, platforms.

\textbf{Report 1: How does money flow?} Periodic tribute without inspecting cultivation, or household registration and direct village taxation? The first is cheap but rough, the second penetrating but expensive.

\textbf{Report 2: How are people appointed?} Transferable central officials, retained local princes, or garrisons? Outsiders strengthen control at staffing cost; agents save money but pursue interests of their own.

\textbf{Report 3: How do routes connect?} Roads, relay stations, shared administrative writing carry orders, troops, documents like vessels. Beyond their reach, map colour is mere paint.

\textbf{Report 4: How are rules set?} How far do measures, currency, law, official language converge? Universal uniformity or local ways?

Law especially reveals imperial intent. Earlier than all this chapter's empires, eighteenth-century BCE Hammurabi displayed nearly three hundred provisions on stone, declaring divine election to establish justice and protect weak from strong. Whatever enforcement, stone announces rules independent of place or official mood, binding whoever comes. Empire seeks precisely this: royal will becoming the way things naturally ought to be. Qin extends the route with detailed penalties for late documents and measurement errors. Reprimanded Liye clerks are splashes from that legal net closing.

Give our five empires their reports:

\begin{tabularx}{\linewidth}{llllX}
\toprule
Empire & Money & \shortstack[l]{Appoint-\\ments} & \shortstack[l]{Routes and\\language} & Approach to difference \\
\midrule
\shortstack[l]{Akkad\\(c. 2300\\BCE)} & \shortstack[l]{City\\tribute} & \shortstack[l]{Akkadian\\governors} & \shortstack[l]{Akkadian\\official\\documents} & Reshape through the conqueror's language \\
\shortstack[l]{Egyptian\\New\\Kingdom\\(c. 1550--\\1070 BCE)} & \shortstack[l]{Direct\\Nubian tax;\\northern\\tribute} & \shortstack[l]{Nubian\\viceroy;\\Syrian--\\Palestinian\\princes\\retained} & \shortstack[l]{Garrisons\\and\\travelling\\envoys} & Layered: direct nearby rule, distant overlordship \\
\shortstack[l]{Assyria\\(c. 900--\\600 BCE)} & \shortstack[l]{Provincial\\taxes and\\tribute} & \shortstack[l]{Royal-\\appointed\\governors} & \shortstack[l]{Relay roads;\\deportation\\disrupts\\identity} & Level differences through uprooting \\
% EN-PAGINATION-BEGIN empire-comparison
\bottomrule
\end{tabularx}

\begin{tabularx}{\linewidth}{llllX}
\toprule
Empire & Money & \shortstack[l]{Appoint-\\ments} & \shortstack[l]{Routes and\\language} & Approach to difference \\
\midrule
% EN-PAGINATION-END empire-comparison
\shortstack[l]{Persia\\(c. 550--\\330 BCE)} & \shortstack[l]{Fixed\\provincial\\taxes under\\Darius} & \shortstack[l]{Satraps and\\royal\\inspectors,\\the king's\\eyes} & \shortstack[l]{Royal Road\\over 2,000 km;\\relays;\\Aramaic;\\gold darics} & Retain differences to serve empire \\
\shortstack[l]{Qin--Han\\(from 221\\BCE)} & \shortstack[l]{Registered\\households;\\household\\and land tax} & \shortstack[l]{Non-\\hereditary\\central\\appoint-\\ments to\\comman-\\deries and\\counties} & \shortstack[l]{Imperial\\roads;\\common\\script, axle\\gauge,\\measures} & Penetrate every household's daily routines \\
\bottomrule
\end{tabularx}
Study the table for a minute.

First, over a thousand years from Akkad to Qin, answers broadly grow deeper and finer. Not inevitable progress: penetration raises costs, extraction, constraints. Second, no answer is absolutely correct. Persian retention saved resources and maintained two centuries of coexistence, but powerful locals escaped central reach; Qin's efficient penetration grew so tense it collapsed under its second emperor. Third, failures appear diagnostically before maps change: stalled revenue, secessionist agents, broken roads, rejected rules.

Move the tool forward. Examine school funding, teacher and year-head appointment, notices and grade circulation, uniform schedules and phone rules. Then delivery platforms: riders and merchants resemble registered county households; algorithms resemble new measures. Managing difference never left humanity's examination paper.

\section[Evidence: An Order and a Cylinder]{A Little Evidence: An Order and a Cylinder}
Start primary material with Qin. Sima Qian's "Basic Annals of the First Emperor of Qin" records unification in fourteen dry Chinese characters:

\begin{quote}
Unify standards, weights, and measures; give carts one gauge and writing one script.
\end{quote}
Three measures: common length and weight standards, with shi and zhang/chi described as capacity and length units; common axle spacing so carts follow shared ruts without catching or overturning on muddy roads; prescribed official and book script replacing six states' variant forms.

Notice the list form: no motives, exclamations, merely an administrative contents page. Tone itself evidences Qin's vision of unification as \textbf{procedures}, altering specifications item by item until the world aligns. Sima Qian stood only decades after the emperor; his unembellished account exposes a worldview treating difference as a technical problem solved by specifications.

Now another tone. Around the same period in Mesopotamia, Cyrus ordered an inscription buried in Babylon's wall foundations after conquest, today's museum cylinder. Its sense: Cyrus, king of all lands, divinely chosen, entered without bloodshed, restored displaced gods to temples and scattered peoples to homes.

Opposite tones: Qin lists specifications without gods or emotion; Cyrus fills restoration with divine grace. Not cruelty versus kindness, but two \textbf{self-justifications}: making everything uniform, or restoring everything's proper state. One repairs forward, one backward, both explaining why they should rule.

Even unity changes meaning. Qin's qi means alignment; Cyrus's an means peace, everything in place under his authority. Before asking truth, ask \textbf{what the source wants its audience to believe}. Lists make unity procedural; cylinders make conquest gracious. A source need not lie to be limited by standpoint. Neither sees Qianling's clerk or riverside exiles.

\section{The Same Roads, Three Accounts}
Our empires independently built roads, relays, scripts, and unified or tolerated measures and money. Facts have durable evidence: Persian road distances, Liye slips, Behistun cliffs. Disputes concern \textbf{why these actions recur} and \textbf{what they meant to contemporaries}. Compare three accounts with strengths and limits.

\textbf{Explanation 1: Revenue and troops: function.}

Roads accelerate armies and grain; relays outrun rebellions with intelligence and orders; standards prevent tax shrinkage and endless procurement conversions; script prevents distorted frontier commands. Every unification yields fiscal or military returns. Clear and concrete, it explains projects following conquest. But why costly symbolic self-promotion, Darius's cliff or Qin's commemorative stelae, which collect no extra measure of tax? If all is practical, every empire's monuments seem redundant.

\textbf{Explanation 2: Fulfil a universal order: worldview.}

Rulers may sincerely hold cosmic missions: divine Mesopotamian agents must pacify chaos; Chinese Sons of Heaven receive responsibility for all directions, treating division and divergent customs as disorder rather than options. Roads, scripts, inscriptions become \textbf{fulfilling office}, reporting completed world-order work to heaven. This explains non-revenue symbols, but worldview may decorate prior ambitions afterward. We cannot enter ancient minds to separate belief from rhetoric.

\textbf{Explanation 3: One project means different things in different positions.}

Frontier perspectives ask on whom the road is built. Common script helps traders but invalidates a Chu historian's lifelong craft. Registration gives peace-seekers security but labouring households an inescapable net. Herodotus admires Persian couriers undeterred by rain, snow, heat, night; villages supplying horses, grain, labour may see something else. This prevents packaging policies as simply good or bad. Overused, however, it translates everything into oppression and misses benefits: Qianling's new-script document can reach Xianyang, unprecedented connection.

Function sees accounts, worldview banners, position those beneath both. Not indecision but method: \textbf{explaining institutions does not defend them; understanding a ruler's worldview does not require agreement}. Keep all three present.

\section[Further Challenge: Walls or Doorways?]{A Deeper Challenge: Shouting from Walls, or Entering Every Door?}
Our strongest tool is Michael Mann's distinction between two kinds of state power.

\textbf{Despotic power} lets rulers command without consultation: move a village, execute a person. Like shouting from walls, its voice travels widely but afterward people return home to daily routines.

\textbf{Infrastructural power} embeds authority in registers, measures, roads, shared writing, stamped documents. It need not shout: every form, grain measure, official road already routes life through its formats.

Applied to our empires, Assyria towers in despotic force, deporting, besieging, intimidating, but with limited everyday penetration after people move. Qin excels infrastructurally: registers encompass individuals, appointed counties transmit orders, standards enter each inscription, journey, transaction. Liye's reprimanded clerk fingerprints this power; the centre knows a particular minor mistake's date. Persia takes a middle course: coercion governs satraps and revolts, infrastructure builds roads, fixes taxes, issues gold, then stops before every village's daily life.

A counterintuitive warning follows. Embedded power sounds gentler than violent shouting: merely registration and stamps. Yet visible despotism can be located, avoided, denounced, mourned beside rivers. A complete infrastructural net makes even the opponent unclear. Qianling's clerk suffers no direct persecution, merely discovers that without new script, measures, forms he cannot find work. Conquerors' faces disappear; formats remain. Hence Mann's continuing relevance: \textbf{modern state strength lies not merely in fighting, but in registering}. Birth certificates, school records, travel documents automatically archive each step far beyond Qin's dreams.

Do not fear forms as such. Without infrastructure, no compulsory education, emergency grain, vaccination network. Standards and roads also deliver public services. The harder task is to ask when convenience, unity, your welfare are invoked: a voice on the wall, or boxes inside the door? Where should the boxes stop? From Qianling onward, no standard answer exists.

\section[Today: Empires in Your Pocket]{Back to Today: How Many Empires Are in Your Pocket?}
The ancient question operates daily.

\textbf{Mandarin and dialect.} Common speech lets northeastern and Cantonese speakers telephone and use shared textbooks, a genuinely helpful modern common script. Yet dialects rapidly fade and some children cannot understand grandmothers' stories. Unity and loss remain one coin, its measures now sounds.

\textbf{Phone standards.} Converging chargers, shared networks, QR codes arise from engineers repeatedly fixing global axle gauges. Platforms resemble Qin rules: identical algorithms time, price, penalise millions of riders and traders, registering households down to remaining delivery seconds. Persian arrangements exist too, letting local agents run sectors in exchange for commissions. Our four reports still apply.

\textbf{Europe's experiment.} The European Union shares a euro across nineteen countries while retaining twenty-four official languages and national laws, textbooks, anthems, a monetary unity Persia might appreciate. Debates over common budgets, refugee quotas, borders each match a diagnostic column. Unfinished, the ancient exam continues.

\textbf{Your school.} Shared schedules, exams, uniforms serve fairness and efficiency; clubs, dialects, hometown food, friendship circles protect differences. Each blanket rule and exception miniaturises imperial questions.

From Liye wood to QR codes, tools change, questions remain: which things must converge for mass order, and which uniformities kill languages, crafts, memories, human origins?

\section{Your Turn: Where Would You Draw the Line?}
Return to the wooden slip.

The nameless clerk leaves his question: designing a new school's rules, online community, or multiethnic city's administration, where would you require unity and where protect difference?

Count your cards. Standards support fair trade, exams, justice; ending division's violence cost countless lives; shared systems carry distant public resources. Yet each erased difference may conceal a weeping exile; centres define while frontiers pay; Persian experiments show preserved difference can be intelligent order, not chaos.

Harder still, positions calculate differently. Couriers' successful roads burden villages; deepest power often quietly lays boxes beneath your feet rather than shouting.

Unify language? Measures? Law? Faith, food, holidays, weddings, funerals? Who pays for each wrongly drawn line?

No standard answer. Serious answering does something Qianling's clerk could not: stand in every position and reckon everyone's account. That is why the humanities exist.
